\documentclass[11pt,a4paper]{article}

\usepackage{amsmath}
\usepackage{amssymb}

\usepackage[spanish,es-tabla]{babel}
\usepackage[utf8]{inputenc}
\usepackage[T1]{fontenc}

\usepackage{graphicx}
\usepackage{subcaption}
\usepackage[font=footnotesize,labelfont=small]{caption}
\usepackage{float}
\usepackage{placeins}
\usepackage{multirow}
\usepackage[table]{xcolor}
\usepackage[export]{adjustbox}

\usepackage{multicol}
\usepackage[a4paper, left=20mm, right=20mm, top=10mm, bottom=30mm]{geometry}

\usepackage{verbatim}
\usepackage{appendix}
\usepackage{hyperref}

\usepackage{csquotes}
\usepackage[backend=biber,style=numeric-comp,sorting=none]{biblatex}
\addbibresource{references.bib}

%\usepackage{svg}
%\usepackage{circuitikz}
%\usepackage{commath}
%\usepackage{blindtext}
%\usepackage{lipsum}
%\usepackage{notoccite}

\begin{document}

\begin{center}		
\vspace{1cm}
\textsc{\LARGE{Construcción de polarímetro dinámico\\}}
\vspace{.5cm}
{ \large \bfseries Grupo 13\\}
\vspace{0.2cm}
{ \large \bfseries Barros Martina, Pironio Bernardo\\}
\vspace{0.2cm}
{ \large \textit{ bpironio@gmail.com, barrosmartina.bm@gmail.com}\\}
\vspace{0.2cm}
{\large Laboratorio de iones y atomos frios, $1^{er}$ cuatrimestre 2026\\
Departamento de Física, FCEyN, UBA\\}
\vspace{0.2cm}
\today
\vspace{0.5cm}

\begin{abstract}
\end{abstract}
\end{center}

\end{document}